\begin{helper}
Let $\mathcal F$ be a finite multi-hypergraph with distinct edge labels
$f$ and $g$.  Suppose $u$ and $v$ are distinct old vertices, with
$u\in f\cap g$ and $v\in g$.  Form the split-edge hypergraph
\[
\mathcal F_2=\noderef{KUniformSplitEdgeHypergraph}(k,\mathcal F,g,v,x_1).
\]
Then the independent-triple count in the line-graph neighborhood of $f$
does not decrease after passing to the old-edge copy of $f$:
\[
T_{L(\mathcal F)}(f)
\le
T_{L(\mathcal F_2)}(\operatorname{inl} f).
\]
\end{helper}
\begin{proof}
By \noderef{IndependentTripleCount}, the left side counts three-element
subsets of the line-graph neighborhood of $f$ whose members are pairwise
nonadjacent in \(L(\mathcal F)\).  Send such a triple \(S\) to the triple
\(\operatorname{inl}(S)=\{\operatorname{inl}(e):e\in S\}\) of old-edge
copies in \(\mathcal F_2\).

The map preserves size three because the old-edge-label inclusion
\(e\mapsto \operatorname{inl}(e)\) is injective.  It also sends every old
neighbor of \(f\) to a neighbor of \(\operatorname{inl}(f)\).  Indeed, by
\noderef{LineGraphOfHypergraph}, adjacency means distinct labels with
intersecting hyperedges.  If the old neighbor is not \(g\), then
\noderef{KUniformSplitEdgeHypergraph} says both hyperedges are simply copied
on old vertices, so their nonempty intersection is copied.  If the old
neighbor is \(g\), then the modified copy of \(g\) deletes \(v\) and adds the
fresh vertex \(x_1\); the vertex \(u\) is still present because \(u\in g\)
and \(u\ne v\), and it also lies in \(f\).  Thus \(\operatorname{inl}(g)\)
still meets \(\operatorname{inl}(f)\).

It remains to check pairwise independence.  For any two old labels \(a,b\),
if their old-edge copies are adjacent after the split, then their split
hyperedges have a common vertex.  When neither label is \(g\), this common
vertex is the copy of an old vertex lying in both original hyperedges.  When
one label is \(g\), the other copied old hyperedge contains no fresh vertex,
so any common vertex again comes from an old vertex of \(g\) that was not
deleted and lies in the other original hyperedge.  When both labels are
\(g\), the original hyperedge \(g\) is nonempty because it contains \(v\).
In every case, adjacency of old-edge copies in \(L(\mathcal F_2)\) implies
adjacency of the original labels in \(L(\mathcal F)\).  Therefore a pairwise
nonadjacent old triple maps to a pairwise nonadjacent split triple.

Finally, the map on triples is injective because old-edge labels are copied
injectively.  Hence the finite set counted by
\(T_{L(\mathcal F)}(f)\) injects into the finite set counted by
\(T_{L(\mathcal F_2)}(\operatorname{inl} f)\), proving the stated
inequality.
\end{proof}
